% English abstract (draft of 2026-11-17, realigned on 2026-11-19 with Chapters 5 and 6), a faithful translation of
% content/summary.tex.
Event cameras record brightness changes at the moment they occur, producing a signal whose information lies largely in the temporal order of activity. Spiking Transformers are natural candidates to process it, yet they handle time mostly implicitly, and making the past explicitly accessible costs memory, a scarce resource in compact models meant for FPGAs. This thesis asks how a budget of about half a million parameters should be spent so that a Spiking Transformer captures the temporal relations that matter in event data, while keeping the cost of memory explicit.

The work changes a single factor at a time, on four neuromorphic datasets. Reallocating the budget of a compact QKFormer raises the Macro-F1 on DVS-Lip from 44 to 55 points with only 320 additional parameters, and a variable-length input, in which every step has the same physical duration, gives the network's delays a precise temporal meaning. Building on this, the thesis proposes an interval relational memory, the antisymmetric interval-relational mixer (IRM), which summarizes the past in causal bands at fixed physical distances and compares it with the present through a low-rank bilinear term inspired by the motion correlator of Hassenstein and Reichardt. On THU-EACT-50-CHL the IRM improves the reference model by 7.6 Macro-F1 points on validation and by 10.8 on the official test. The gain comes from the comparison rather than from the added parameters, and on the test set it remains positive with one and two million parameters, though smaller, as expected from a module designed mainly for the compact regime. Once compressed after training, the IRM costs less than the mixer it replaces at every width. With one to two million parameters it outperforms on the official test all published methods without external data (74.3\% against 69.9\%), and with 0.46 million it exceeds the best known spiking network by almost nineteen points. On DVS-Lip, where the compact model is limited by the variance of the solution rather than by memory, distillation adds up with a comparison of the same kind computed before pooling, between neighbouring cells and over a few tens of milliseconds, and together they raise the compact model from 50.8\% to 56.8\% on the test set, while with two million parameters the network reaches 60.2\%.

With costs estimated rather than measured on hardware, the results indicate that in a compact Spiking Transformer temporal capability depends less on the amount of memory than on how the past is compared with the present, once each step has been given a precise physical meaning. The same oriented comparison helps both tasks, but at the scale each one requires, over distances of seconds for actions and between neighbouring details, over a few tens of milliseconds, for lip reading.
